\section{蜀汉的 dossier：盆地、山口与一支小国的选择}
\label{sec:shu-expeditions}

\begin{event}{\eventanchor{TK-0221-SHU-HAN-FOUNDATION}{刘备成都称帝}}{章武元年四月（二二一年）}{蜀汉 · 成都}{即位年月可核，礼仪与承认范围有争议}
《三国志·先主传》《三国志·后主传》与《华阳国志·刘先主志》支持成都即位和建元章武的基本年序。国号礼仪的细部以及这一名义在益州、汉中之外获得何种实际承认，不能由称帝本身预先回答。
\end{event}

\subsection{纲要：名义必须由后方供养}

\hyperlink{evt:TK-0221-SHU-HAN-FOUNDATION}{刘备在成都称帝}给蜀汉以汉统语言，却
没有改变它的资源尺度。夷陵后的永安继承和
\eventanchor{TK-0223-SHU-WU-RECONCILIATION}{蜀吴复通}，迫使诸葛亮先把外交、南中
和成都的行政重新排出次序。二二五年的
\eventanchor{TK-0225-SHU-NANZHONG}{南征}可以理解为这种重新排序：它所追求的
不是将南中化作没有差异的内地，而是让盆地国家在地方合作、交通和征发上重新
获得余地。

汉中是蜀汉的前沿，也是它的负担。诸葛亮驻汉中的准备、祁山的机会与街亭后的
撤回，连成\eventanchor{TK-0227-ZHUGE-LIANG-HANZHONG}{组织北伐}到
\eventanchor{TK-0228-JIETING}{承认补给限度}的过程。取武都、阴平固然增加屏障，
但不能把山道变短；二三一年李严因转输问题被废，正说明军事决策要通过地方、
仓道和文书才会落地。

\subsection{留守、节制与辅政的断裂}

建兴五年，诸葛亮令\eventanchor{TK-DENS-0227-01}{李严留守永安，自己驻军汉中}。
这不是把蜀汉简单分成“北伐派”与“守成派”：永安要照看峡口和江防，汉中则要
把人粮送向山口，两处责任必须同时有人承担。纪传可以确定分置，不能据此补写
两人对战略的全部意见；它至少显示成都无法把盆地外缘当作一条单线。

五丈原以后，费祎在成都主持政务，并在延熙十一年
\eventanchor{TK-DENS-0248-01}{限制大规模北伐而维持汉中准备}。这是一种分配
风险的选择，而不是前线从此静止：守备、道路和边郡关系仍要花费资源。延熙十六年，
这位辅政者为魏降人郭脩所杀，\eventanchor{TK-DENS-0253-01}{成都失去一名能够
协调中枢与前线的人}；记载所能确定的是刺杀与继任的顺序，不能由此推知此后
所有北伐政策都出自同一原因。

\subsection{北伐不是一种恒定政策}

五丈原的\eventanchor{TK-0234-WUZHANG}{相持与撤军}后，蒋琬接掌中枢，北伐并未
从国家议程中消失，然而其节奏和能承受的成本变得不同。姜维二四七年借陇西
局势出兵，二六二年又在狄道受挫；这些行动既利用边缘集团带来的缝隙，也持续
暴露成都至陇右的转输压力。将其全归结为某位将领的好战或怯战，会遮蔽小国
必须以有限后方维持前线的难题。

二六三年的魏军多路南下终于使这种约束集中爆发。
\eventanchor{TK-0263-CHENGDU-SURRENDER}{刘禅出降}结束蜀汉政权，却不是蜀地社会
立即被抹平；随后钟会之乱已说明接收者也须学习如何使用道路、降众与地方信息。

\subsection{制度得失}

蜀汉以丞相府和汉中前线把盆地资源集中起来，能够在强邻之间保持相当长的
行动能力；代价是前线每一次延伸都加重道路和后方的负担。它的经验不在于北伐
是否应当被道德化评断，而在于政治名义、边郡关系和转输能力必须同时成立。

\begin{sourcenote}[南中与山地材料]
蜀汉年序以《三国志·蜀书》为主，南中和地方关系还须以《华阳国志》及区域
研究校读。地方叙事可以扩展问题，但不能据传说补写行军细节，或把一次军政
重置等同于长期、完全的服从。
\end{sourcenote}


\subsection{蜀汉纪年补链：盆地朝廷与山道前线之间}

以下蜀汉补链将成都朝廷的处置、山道前线与边郡接触接回盆地后方的限制。锚点只保留本纪可见的顺序，不把其沉默补成郡县服从、行军细节或将领动机。

\paragraph{章武三年四月（223年）} 《本纪》在章武三年四月（223年）先把这些动作排入同一条年序：\eventanchor{TK-0223-LIU-BEI-DEATH}{刘备在永安去世，刘禅继位，诸葛亮受遗诏辅政}。\footnote{核心来源：《三国志·先主传》；《三国志·后主传》；《三国志·诸葛亮传》。材料限度：卒年、继位和辅政事实可靠；遗诏措辞、永安部署及刘备病因有争议}

\paragraph{建兴六年春（228年）} 到了建兴六年春（228年），朝廷可见的变化并不只在一项大事上，\eventanchor{TK-0228-FIRST-NORTHERN-EXPEDITION}{诸葛亮出祁山，南安、天水、安定一度响应蜀汉，魏调张郃等西援}。\footnote{核心来源：《三国志·诸葛亮传》；《三国志·明帝纪》；《三国志·张郃传》；《资治通鉴》魏纪。材料限度：出兵与数郡震动可靠；响应范围、地方倒向的原因和双方兵力有争议}

\paragraph{建兴七年（229年）} 记录写下建兴七年（229年）的多个接口：\eventanchor{TK-0229-WUDU-YINPING}{诸葛亮取武都、阴平二郡，魏将郭淮退却，蜀汉扩展汉中西北屏障}。\footnote{核心来源：《三国志·诸葛亮传》；《三国志·郭淮传》；《资治通鉴》魏纪。材料限度：取郡结果大体可靠；郡界、守令转换和郭淮撤军原因有争议}

\paragraph{建兴九年（231年）} 建兴九年（231年）的条目把命令、任免或接触并列起来：\eventanchor{TK-0231-QISHAN-EXPEDITION}{诸葛亮再次出祁山，与司马懿部相持，因粮运不继而撤军}，并\eventanchor{TK-0231-LI-YAN-DISMISSAL}{李严因粮运失职并伪称奉诏召军，被诸葛亮弹劾废为平民}。\footnote{核心来源：《三国志·诸葛亮传》；《三国志·司马懿传》；《三国志·张郃传》；《资治通鉴》魏纪；《三国志·李严传》；《三国志·诸葛亮传》；《三国志·李严传》裴松之注引《诸葛亮集》。材料限度：出兵、对峙和撤军可靠；木牛流马、卤城战果和双方伤亡有争议}

\paragraph{建兴十二年（234年）} 沿着建兴十二年（234年）留下的纪年节点，可以看到资源和权力怎样被逐项调度：\eventanchor{TK-0234-JIANG-WAN-REGENT}{诸葛亮死后，蒋琬录尚书事，继掌蜀汉中枢政务}。\footnote{核心来源：《三国志·蒋琬传》；《三国志·后主传》；《华阳国志·刘后主志》。材料限度：接掌中枢事实可靠；蒋琬权力起点、与费祎和姜维的分工及政策连续性有争议}

\paragraph{延熙十年（247年）} 同年的记载没有替行动者说明动机，却留下了可接续的顺序：\eventanchor{TK-0247-JIANG-WEI-LONGXI}{姜维趁陇西羌胡起事出兵，魏将郭淮等应战，蜀汉重新将北伐重心推向陇右}。\footnote{核心来源：《三国志·姜维传》；《三国志·郭淮传》；《资治通鉴》魏纪。材料限度：出兵和魏蜀交战可靠；羌胡各部关系、地方响应和姜维战略自主性有争议}

\paragraph{景耀五年（262年）} 《本纪》在景耀五年（262年）先把这些动作排入同一条年序：\eventanchor{TK-0262-JIANG-WEI-DIDAO}{姜维围攻狄道，邓艾据险救援，蜀军因粮运和魏军反击撤退}。\footnote{核心来源：《三国志·姜维传》；《三国志·邓艾传》；《三国志·后主传》；《资治通鉴》魏纪。材料限度：围城、邓艾救援和撤军可靠；兵力、战果、郭淮角色与撤军损失有争议}

\paragraph{章武元年四月（221年）} 到了章武元年四月（221年），朝廷可见的变化并不只在一项大事上，\eventanchor{TK-LEDGER-002-1}{“刘备在成都即皇帝位，建元章武，以汉室继承者自居”的前期动员与资源准备}；\eventanchor{TK-LEDGER-002-2}{“刘备在成都即皇帝位，建元章武，以汉室继承者自居”的命令、联盟或地方响应}；\eventanchor{TK-LEDGER-002-3}{刘备在成都即皇帝位，建元章武，以汉室继承者自居}；\eventanchor{TK-LEDGER-002-4}{“刘备在成都即皇帝位，建元章武，以汉室继承者自居”的执行中的空间与人员处置}；\eventanchor{TK-LEDGER-002-5}{“刘备在成都即皇帝位，建元章武，以汉室继承者自居”的后续制度或军政调整}，并\eventanchor{TK-LEDGER-002-6}{“刘备在成都即皇帝位，建元章武，以汉室继承者自居”的异源材料与影响范围的复核}。\footnote{核心来源：《三国志·先主传》；《三国志·后主传》；《华阳国志·刘先主志》。材料限度：核心事项已有既有账本来源；本分解节点的参与者、执行范围和时间先后仍须逐案核验}

\paragraph{章武三年四月（223年）} 记录写下章武三年四月（223年）的多个接口：\eventanchor{TK-LEDGER-006-1}{“刘备在永安去世，刘禅继位，诸葛亮受遗诏辅政”的前期动员与资源准备}；\eventanchor{TK-LEDGER-006-2}{“刘备在永安去世，刘禅继位，诸葛亮受遗诏辅政”的命令、联盟或地方响应}；\eventanchor{TK-LEDGER-006-3}{刘备在永安去世，刘禅继位，诸葛亮受遗诏辅政}；\eventanchor{TK-LEDGER-006-4}{“刘备在永安去世，刘禅继位，诸葛亮受遗诏辅政”的执行中的空间与人员处置}；\eventanchor{TK-LEDGER-006-5}{“刘备在永安去世，刘禅继位，诸葛亮受遗诏辅政”的后续制度或军政调整}，并\eventanchor{TK-LEDGER-006-6}{“刘备在永安去世，刘禅继位，诸葛亮受遗诏辅政”的异源材料与影响范围的复核}。\footnote{核心来源：《三国志·先主传》；《三国志·后主传》；《三国志·诸葛亮传》。材料限度：核心事项已有既有账本来源；本分解节点的参与者、执行范围和时间先后仍须逐案核验}

\paragraph{建兴三年（225年）} 建兴三年（225年）的条目把命令、任免或接触并列起来：\eventanchor{TK-LEDGER-009-1}{“诸葛亮南征，平定益州郡、永昌等地叛乱并重置地方统治关系”的前期动员与资源准备}；\eventanchor{TK-LEDGER-009-2}{“诸葛亮南征，平定益州郡、永昌等地叛乱并重置地方统治关系”的命令、联盟或地方响应}；\eventanchor{TK-LEDGER-009-3}{诸葛亮南征，平定益州郡、永昌等地叛乱并重置地方统治关系}；\eventanchor{TK-LEDGER-009-4}{“诸葛亮南征，平定益州郡、永昌等地叛乱并重置地方统治关系”的执行中的空间与人员处置}；\eventanchor{TK-LEDGER-009-5}{“诸葛亮南征，平定益州郡、永昌等地叛乱并重置地方统治关系”的后续制度或军政调整}，并\eventanchor{TK-LEDGER-009-6}{“诸葛亮南征，平定益州郡、永昌等地叛乱并重置地方统治关系”的异源材料与影响范围的复核}。\footnote{核心来源：《三国志·诸葛亮传》；《三国志·马谡传》裴松之注引《襄阳记》；《华阳国志·南中志》。材料限度：核心事项已有既有账本来源；本分解节点的参与者、执行范围和时间先后仍须逐案核验}

\paragraph{建兴五年（227年）} 沿着建兴五年（227年）留下的纪年节点，可以看到资源和权力怎样被逐项调度：\eventanchor{TK-LEDGER-011-1}{“诸葛亮驻汉中并上《出师表》，以北伐准备重新组织蜀汉军政”的前期动员与资源准备}；\eventanchor{TK-LEDGER-011-2}{“诸葛亮驻汉中并上《出师表》，以北伐准备重新组织蜀汉军政”的命令、联盟或地方响应}；\eventanchor{TK-LEDGER-011-3}{诸葛亮驻汉中并上《出师表》，以北伐准备重新组织蜀汉军政}；\eventanchor{TK-LEDGER-011-4}{“诸葛亮驻汉中并上《出师表》，以北伐准备重新组织蜀汉军政”的执行中的空间与人员处置}；\eventanchor{TK-LEDGER-011-5}{“诸葛亮驻汉中并上《出师表》，以北伐准备重新组织蜀汉军政”的后续制度或军政调整}，并\eventanchor{TK-LEDGER-011-6}{“诸葛亮驻汉中并上《出师表》，以北伐准备重新组织蜀汉军政”的异源材料与影响范围的复核}。\footnote{核心来源：《三国志·诸葛亮传》；《三国志·后主传》；《出师表》传本。材料限度：核心事项已有既有账本来源；本分解节点的参与者、执行范围和时间先后仍须逐案核验}

\paragraph{建兴六年春（228年）} 同年的记载没有替行动者说明动机，却留下了可接续的顺序：\eventanchor{TK-LEDGER-012-1}{“诸葛亮出祁山，南安、天水、安定一度响应蜀汉，魏调张郃等西援”的前期动员与资源准备}；\eventanchor{TK-LEDGER-012-2}{“诸葛亮出祁山，南安、天水、安定一度响应蜀汉，魏调张郃等西援”的命令、联盟或地方响应}；\eventanchor{TK-LEDGER-012-3}{诸葛亮出祁山，南安、天水、安定一度响应蜀汉，魏调张郃等西援}；\eventanchor{TK-LEDGER-012-4}{“诸葛亮出祁山，南安、天水、安定一度响应蜀汉，魏调张郃等西援”的执行中的空间与人员处置}；\eventanchor{TK-LEDGER-012-5}{“诸葛亮出祁山，南安、天水、安定一度响应蜀汉，魏调张郃等西援”的后续制度或军政调整}，并\eventanchor{TK-LEDGER-012-6}{“诸葛亮出祁山，南安、天水、安定一度响应蜀汉，魏调张郃等西援”的异源材料与影响范围的复核}。\footnote{核心来源：《三国志·诸葛亮传》；《三国志·明帝纪》；《三国志·张郃传》；《资治通鉴》魏纪。材料限度：核心事项已有既有账本来源；本分解节点的参与者、执行范围和时间先后仍须逐案核验}

\paragraph{建兴六年（228年）} 《本纪》在建兴六年（228年）先把这些动作排入同一条年序：\eventanchor{TK-LEDGER-013-1}{“马谡失守街亭，诸葛亮撤回汉中并处置失军者”的前期动员与资源准备}；\eventanchor{TK-LEDGER-013-2}{“马谡失守街亭，诸葛亮撤回汉中并处置失军者”的命令、联盟或地方响应}；\eventanchor{TK-LEDGER-013-3}{马谡失守街亭，诸葛亮撤回汉中并处置失军者}；\eventanchor{TK-LEDGER-013-4}{“马谡失守街亭，诸葛亮撤回汉中并处置失军者”的执行中的空间与人员处置}；\eventanchor{TK-LEDGER-013-5}{“马谡失守街亭，诸葛亮撤回汉中并处置失军者”的后续制度或军政调整}，并\eventanchor{TK-LEDGER-013-6}{“马谡失守街亭，诸葛亮撤回汉中并处置失军者”的异源材料与影响范围的复核}。\footnote{核心来源：《三国志·马谡传》；《三国志·王平传》；《三国志·张郃传》；《三国志·诸葛亮传》裴松之注。材料限度：核心事项已有既有账本来源；本分解节点的参与者、执行范围和时间先后仍须逐案核验}

\paragraph{建兴七年（229年）} 到了建兴七年（229年），朝廷可见的变化并不只在一项大事上，\eventanchor{TK-LEDGER-016-1}{“诸葛亮取武都、阴平二郡，魏将郭淮退却，蜀汉扩展汉中西北屏障”的前期动员与资源准备}；\eventanchor{TK-LEDGER-016-2}{“诸葛亮取武都、阴平二郡，魏将郭淮退却，蜀汉扩展汉中西北屏障”的命令、联盟或地方响应}；\eventanchor{TK-LEDGER-016-3}{诸葛亮取武都、阴平二郡，魏将郭淮退却，蜀汉扩展汉中西北屏障}；\eventanchor{TK-LEDGER-016-4}{“诸葛亮取武都、阴平二郡，魏将郭淮退却，蜀汉扩展汉中西北屏障”的执行中的空间与人员处置}；\eventanchor{TK-LEDGER-016-5}{“诸葛亮取武都、阴平二郡，魏将郭淮退却，蜀汉扩展汉中西北屏障”的后续制度或军政调整}，并\eventanchor{TK-LEDGER-016-6}{“诸葛亮取武都、阴平二郡，魏将郭淮退却，蜀汉扩展汉中西北屏障”的异源材料与影响范围的复核}。\footnote{核心来源：《三国志·诸葛亮传》；《三国志·郭淮传》；《资治通鉴》魏纪。材料限度：核心事项已有既有账本来源；本分解节点的参与者、执行范围和时间先后仍须逐案核验}

\paragraph{建兴九年（231年）} 记录写下建兴九年（231年）的多个接口：\eventanchor{TK-LEDGER-021-1}{“诸葛亮再次出祁山，与司马懿部相持，因粮运不继而撤军”的前期动员与资源准备}；\eventanchor{TK-LEDGER-021-2}{“诸葛亮再次出祁山，与司马懿部相持，因粮运不继而撤军”的命令、联盟或地方响应}；\eventanchor{TK-LEDGER-021-3}{诸葛亮再次出祁山，与司马懿部相持，因粮运不继而撤军}；\eventanchor{TK-LEDGER-021-4}{“诸葛亮再次出祁山，与司马懿部相持，因粮运不继而撤军”的执行中的空间与人员处置}；\eventanchor{TK-LEDGER-021-5}{“诸葛亮再次出祁山，与司马懿部相持，因粮运不继而撤军”的后续制度或军政调整}；\eventanchor{TK-LEDGER-021-6}{“诸葛亮再次出祁山，与司马懿部相持，因粮运不继而撤军”的异源材料与影响范围的复核}；\eventanchor{TK-LEDGER-022-1}{“李严因粮运失职并伪称奉诏召军，被诸葛亮弹劾废为平民”的前期动员与资源准备}；\eventanchor{TK-LEDGER-022-2}{“李严因粮运失职并伪称奉诏召军，被诸葛亮弹劾废为平民”的命令、联盟或地方响应}；\eventanchor{TK-LEDGER-022-3}{李严因粮运失职并伪称奉诏召军，被诸葛亮弹劾废为平民}；\eventanchor{TK-LEDGER-022-4}{“李严因粮运失职并伪称奉诏召军，被诸葛亮弹劾废为平民”的执行中的空间与人员处置}；\eventanchor{TK-LEDGER-022-5}{“李严因粮运失职并伪称奉诏召军，被诸葛亮弹劾废为平民”的后续制度或军政调整}，并\eventanchor{TK-LEDGER-022-6}{“李严因粮运失职并伪称奉诏召军，被诸葛亮弹劾废为平民”的异源材料与影响范围的复核}。\footnote{核心来源：《三国志·诸葛亮传》；《三国志·司马懿传》；《三国志·张郃传》；《资治通鉴》魏纪；《三国志·李严传》；《三国志·诸葛亮传》；《三国志·李严传》裴松之注引《诸葛亮集》。材料限度：核心事项已有既有账本来源；本分解节点的参与者、执行范围和时间先后仍须逐案核验}

\paragraph{建兴十二年（234年）} 建兴十二年（234年）的条目把命令、任免或接触并列起来：\eventanchor{TK-LEDGER-025-1}{“诸葛亮出斜谷屯五丈原，与司马懿相持后病逝，蜀军撤回汉中”的前期动员与资源准备}；\eventanchor{TK-LEDGER-025-2}{“诸葛亮出斜谷屯五丈原，与司马懿相持后病逝，蜀军撤回汉中”的命令、联盟或地方响应}；\eventanchor{TK-LEDGER-025-3}{诸葛亮出斜谷屯五丈原，与司马懿相持后病逝，蜀军撤回汉中}；\eventanchor{TK-LEDGER-025-4}{“诸葛亮出斜谷屯五丈原，与司马懿相持后病逝，蜀军撤回汉中”的执行中的空间与人员处置}；\eventanchor{TK-LEDGER-025-5}{“诸葛亮出斜谷屯五丈原，与司马懿相持后病逝，蜀军撤回汉中”的后续制度或军政调整}；\eventanchor{TK-LEDGER-025-6}{“诸葛亮出斜谷屯五丈原，与司马懿相持后病逝，蜀军撤回汉中”的异源材料与影响范围的复核}；\eventanchor{TK-LEDGER-027-1}{“诸葛亮死后，蒋琬录尚书事，继掌蜀汉中枢政务”的前期动员与资源准备}；\eventanchor{TK-LEDGER-027-2}{“诸葛亮死后，蒋琬录尚书事，继掌蜀汉中枢政务”的命令、联盟或地方响应}；\eventanchor{TK-LEDGER-027-3}{诸葛亮死后，蒋琬录尚书事，继掌蜀汉中枢政务}；\eventanchor{TK-LEDGER-027-4}{“诸葛亮死后，蒋琬录尚书事，继掌蜀汉中枢政务”的执行中的空间与人员处置}；\eventanchor{TK-LEDGER-027-5}{“诸葛亮死后，蒋琬录尚书事，继掌蜀汉中枢政务”的后续制度或军政调整}，并\eventanchor{TK-LEDGER-027-6}{“诸葛亮死后，蒋琬录尚书事，继掌蜀汉中枢政务”的异源材料与影响范围的复核}。\footnote{核心来源：《三国志·诸葛亮传》；《三国志·司马懿传》；《三国志·杨仪传》裴松之注引《魏氏春秋》；《三国志·蒋琬传》；《三国志·后主传》；《华阳国志·刘后主志》。材料限度：核心事项已有既有账本来源；本分解节点的参与者、执行范围和时间先后仍须逐案核验}


\subsection{蜀纪所见：成都、汉中与边郡的短条}

蜀书帝纪留下的短条多与山道国家的征发、任命和边郡关系相连。它们让行动的节点可见，不能替代南中、陇右或普通人群自身的记录。

\paragraph{221年} 《三国志·蜀书·先主传》；《三国志·蜀书·后主传》在这一年并列留下官署、任命与诏令的记录。\footnote{核心来源：《三国志·蜀书·先主传》；《三国志·蜀书·后主传》。材料限度：本纪可核对行动的基本顺序；兵额、路线、战损和具体指挥过程仍须另证。；本纪以本政权视角记录对手或边缘群体，不足以断定对方内部决策、疆界或长期归属。；本纪选择朝廷可见事项；不得由一条记录推定各郡社会状况或具体群体动机。；本纪可证实所载命令、任免或结果；不据此推定各郡执行、全部参与者或私人动机。}
\begin{itemize}[leftmargin=1.8em,itemsep=2pt,topsep=3pt]
  \item \eventanchor{TK-SRC-0221-17}{車騎將軍張飛為其左右所害}。
  \item \eventanchor{TK-SRC-0221-18}{初，先主忿孫權之襲關羽，將東征，秋七月，遂帥諸軍伐吳}。
  \item \eventanchor{TK-SRC-0221-19}{今以禪為皇太子，以承宗廟，祗肅社稷}。
  \item \eventanchor{TK-SRC-0221-20}{六月，以子永為魯王，理為梁王}。
  \item \eventanchor{TK-SRC-0221-21}{使使持節丞相亮授印綬，敬聽師傅，行一物而三善皆得焉，可不勉與！」三年夏四月，先主殂於永安宮}。
  \item \eventanchor{TK-SRC-0221-22}{孫權遣書請和，先主盛怒不許，吳將陸議、李異、劉阿等屯巫、秭歸；將軍吳班、馮習自巫攻破異等，軍次秭歸，武陵五谿蠻夷遣使請兵}。
  \item \eventanchor{TK-SRC-0221-23}{五月，後主襲位於成都，時年十七，尊皇后曰皇太后}。
  \item \eventanchor{TK-SRC-0221-24}{五月，立皇后吳氏，子禪為皇太子}。
  \item \eventanchor{TK-SRC-0221-25}{以諸葛亮為丞相，許靖為司徒}。
  \item \eventanchor{TK-SRC-0221-26}{章武元年夏四月，大赦，改年}。
  \item \eventanchor{TK-SRC-0221-27}{置百官，立宗廟，祫祭高皇帝以下}。
\end{itemize}

\paragraph{222年} 《三国志·蜀书·先主传》在这一年并列留下官署、任命与诏令的记录。\footnote{核心来源：《三国志·蜀书·先主传》。材料限度：本纪选择朝廷可见事项；不得由一条记录推定各郡社会状况或具体群体动机。；本纪以本政权视角记录对手或边缘群体，不足以断定对方内部决策、疆界或长期归属。；本纪可证实所载命令、任免或结果；不据此推定各郡执行、全部参与者或私人动机。}
\begin{itemize}[leftmargin=1.8em,itemsep=2pt,topsep=3pt]
  \item \eventanchor{TK-SRC-0222-32}{冬十二月，漢嘉太守黃元聞先主疾不豫，舉兵拒守}。
  \item \eventanchor{TK-SRC-0222-33}{冬十月，詔丞相亮營南北郊於成都}。
  \item \eventanchor{TK-SRC-0222-34}{二年春正月，先主軍還秭歸，將軍吳班、陳式水軍屯夷陵，夾江東西岸}。
  \item \eventanchor{TK-SRC-0222-35}{二月，先主自秭歸率諸將進軍，緣山截嶺，於夷道猇亭駐營，自佷山通武陵，遣侍中馬良安慰五谿蠻夷，咸相率響應}。
  \item \eventanchor{TK-SRC-0222-36}{後十餘日，陸議大破先主軍於猇亭，將軍馮習、張南等皆沒}。
  \item \eventanchor{TK-SRC-0222-37}{孫權聞先主住白帝，甚懼，遣使請和}。
  \item \eventanchor{TK-SRC-0222-38}{吳遣將軍李異、劉阿等踵躡先主軍，屯駐南山}。
  \item \eventanchor{TK-SRC-0222-39}{夏六月，黃氣見自秭歸十餘里中，廣數十丈}。
  \item \eventanchor{TK-SRC-0222-40}{先主許之，遣太中大夫宗瑋報命}。
  \item \eventanchor{TK-SRC-0222-41}{先主自猇亭還秭歸，收合離散兵，遂棄船舫，由步道還魚复，改魚复縣曰永安}。
  \item \eventanchor{TK-SRC-0222-42}{鎮北將軍黃權督江北諸軍，與吳軍相拒於夷陵道}。
\end{itemize}

\paragraph{223年} 《三国志·蜀书·先主传》；《三国志·蜀书·后主传》在这一年并列留下官署、任命与诏令的记录。\footnote{核心来源：《三国志·蜀书·先主传》；《三国志·蜀书·后主传》。材料限度：本纪可证实所载命令、任免或结果；不据此推定各郡执行、全部参与者或私人动机。；本纪选择朝廷可见事项；不得由一条记录推定各郡社会状况或具体群体动机。；本纪可核对行动的基本顺序；兵额、路线、战损和具体指挥过程仍须另证。；本纪以本政权视角记录对手或边缘群体，不足以断定对方内部决策、疆界或长期归属。}
\begin{itemize}[leftmargin=1.8em,itemsep=2pt,topsep=3pt]
  \item \eventanchor{TK-SRC-0223-06}{機權幹略，不逮魏武，是以基宇亦狹}。
  \item \eventanchor{TK-SRC-0223-07}{及其舉國託孤於諸葛亮，而心神無貳，誠君臣之至公，古今之盛軌也}。
  \item \eventanchor{TK-SRC-0223-08}{建興元年夏，臧柯太守朱褒擁郡反}。
  \item \eventanchor{TK-SRC-0223-09}{臨終時，呼魯王與語：「吾亡之後，汝兄弟父事丞相，令卿與丞相共事而已}。
  \item \eventanchor{TK-SRC-0223-10}{遣將軍陳曶討元，元軍敗，順流下江，為其親兵所縛，生致成都，斬之}。
  \item \eventanchor{TK-SRC-0223-11}{遣尚書郎鄧芝固好於吳，吳王孫權與蜀和親使聘，是歲通好}。
  \item \eventanchor{TK-SRC-0223-12}{然折而不撓，終不為下者，抑揆彼之量必不容己，非唯競利，且以避害云爾}。
  \item \eventanchor{TK-SRC-0223-13}{三年春二月，丞相亮自成都到永安}。
  \item \eventanchor{TK-SRC-0223-14}{五月，梓宮自永安還成都，諡曰昭烈皇帝}。
  \item \eventanchor{TK-SRC-0223-15}{夏四月癸巳，先主殂於永安宮，時年六十三}。
  \item \eventanchor{TK-SRC-0223-16}{先是，益州郡有大姓雍闓反，流太守張裔於吳，據郡不賓，越巂夷王高定亦背叛}。
  \item \eventanchor{TK-SRC-0223-17}{先主病篤，託孤於丞相亮，尚書令李嚴為副}。
\end{itemize}

\paragraph{224年} 《三国志·蜀书·后主传》在这一年并列留下赋役、赈济与赏赐的记录。\footnote{核心来源：《三国志·蜀书·后主传》。材料限度：本纪所载多为朝廷命令或结果，不足以证明地方执行、总量和受影响户数。}
\begin{itemize}[leftmargin=1.8em,itemsep=2pt,topsep=3pt]
  \item \eventanchor{TK-SRC-0224-08}{二年春，務農殖穀，閉關息民}。
\end{itemize}

\paragraph{225年} 《三国志·蜀书·后主传》在这一年并列留下官署、任命与诏令的记录。\footnote{核心来源：《三国志·蜀书·后主传》。材料限度：本纪以本政权视角记录对手或边缘群体，不足以断定对方内部决策、疆界或长期归属。}
\begin{itemize}[leftmargin=1.8em,itemsep=2pt,topsep=3pt]
  \item \eventanchor{TK-SRC-0225-06}{改益州郡為建寧郡，分建寧、永昌郡為雲南郡，又分建寧、臧柯為興古郡}。
  \item \eventanchor{TK-SRC-0225-07}{三年春三月，丞相亮南征四郡，四郡皆平}。
\end{itemize}

\paragraph{226年} 《三国志·蜀书·后主传》在这一年并列留下朝廷可见的处置与记述的记录。\footnote{核心来源：《三国志·蜀书·后主传》。材料限度：本纪可核对行动的基本顺序；兵额、路线、战损和具体指挥过程仍须另证。}
\begin{itemize}[leftmargin=1.8em,itemsep=2pt,topsep=3pt]
  \item \eventanchor{TK-SRC-0226-07}{四年春，都護李嚴自永安還住江州，築大城}。
\end{itemize}

\paragraph{227年} 《三国志·蜀书·后主传》在这一年并列留下朝廷可见的处置与记述的记录。\footnote{核心来源：《三国志·蜀书·后主传》。材料限度：本纪以本政权视角记录对手或边缘群体，不足以断定对方内部决策、疆界或长期归属。}
\begin{itemize}[leftmargin=1.8em,itemsep=2pt,topsep=3pt]
  \item \eventanchor{TK-SRC-0227-03}{五年春，丞相亮出屯漢中，營沔北陽平石馬}。
\end{itemize}

\paragraph{228年} 《三国志·蜀书·后主传》在这一年并列留下赋役、赈济与赏赐的记录。\footnote{核心来源：《三国志·蜀书·后主传》。材料限度：本纪所载多为朝廷命令或结果，不足以证明地方执行、总量和受影响户数。；本纪可核对行动的基本顺序；兵额、路线、战损和具体指挥过程仍须另证。；本纪以本政权视角记录对手或边缘群体，不足以断定对方内部决策、疆界或长期归属。}
\begin{itemize}[leftmargin=1.8em,itemsep=2pt,topsep=3pt]
  \item \eventanchor{TK-SRC-0228-07}{冬，復出散關，圍陳倉，糧盡退}。
  \item \eventanchor{TK-SRC-0228-08}{六年春，亮出攻祁山，不克}。
  \item \eventanchor{TK-SRC-0228-09}{魏將王雙率軍追亮，亮與戰，破之，斬雙，還漢中}。
\end{itemize}

\paragraph{229年} 《三国志·蜀书·后主传》在这一年并列留下军队、边地与征讨的记录。\footnote{核心来源：《三国志·蜀书·后主传》。材料限度：本纪以本政权视角记录对手或边缘群体，不足以断定对方内部决策、疆界或长期归属。；本纪可核对行动的基本顺序；兵额、路线、战损和具体指挥过程仍须另证。}
\begin{itemize}[leftmargin=1.8em,itemsep=2pt,topsep=3pt]
  \item \eventanchor{TK-SRC-0229-01}{冬，亮徙從府營於南山下原上，築漢、樂二城}。
  \item \eventanchor{TK-SRC-0229-02}{七年春，亮遣陳式攻武都、陰平，遂克定二郡}。
  \item \eventanchor{TK-SRC-0229-03}{是歲，孫權稱帝，與蜀約盟，共交分天下}。
\end{itemize}

\paragraph{230年} 《三国志·蜀书·后主传》在这一年并列留下赋役、赈济与赏赐的记录。\footnote{核心来源：《三国志·蜀书·后主传》。材料限度：本纪选择朝廷可见事项；不得由一条记录推定各郡社会状况或具体群体动机。；本纪可核对行动的基本顺序；兵额、路线、战损和具体指挥过程仍须另证。；本纪可证实所载命令、任免或结果；不据此推定各郡执行、全部参与者或私人动机。}
\begin{itemize}[leftmargin=1.8em,itemsep=2pt,topsep=3pt]
  \item \eventanchor{TK-SRC-0230-01}{八年秋，魏使司馬懿由西城，張郃由子午，曹真由斜谷，欲攻漢中}。
  \item \eventanchor{TK-SRC-0230-02}{丞相亮待之於城固、赤坂，大雨道絕，真等皆還}。
  \item \eventanchor{TK-SRC-0230-03}{是歲，魏延破魏雍州剌史郭淮於陽溪}。
\end{itemize}

\paragraph{231年} 《三国志·蜀书·后主传》在这一年并列留下官署、任命与诏令的记录。\footnote{核心来源：《三国志·蜀书·后主传》。材料限度：本纪可核对行动的基本顺序；兵额、路线、战损和具体指挥过程仍须另证。；本纪可证实所载命令、任免或结果；不据此推定各郡执行、全部参与者或私人动机。；本纪选择朝廷可见事项；不得由一条记录推定各郡社会状况或具体群体动机。}
\begin{itemize}[leftmargin=1.8em,itemsep=2pt,topsep=3pt]
  \item \eventanchor{TK-SRC-0231-01}{九年春二月，亮復出軍圍祁山，始以木牛運}。
  \item \eventanchor{TK-SRC-0231-02}{秋八月，都護李嚴廢徙梓潼郡}。
  \item \eventanchor{TK-SRC-0231-03}{夏六月，亮糧盡退軍，郃追至青封，與亮交戰，被箭死}。
\end{itemize}

\paragraph{232年} 《三国志·蜀书·后主传》在这一年并列留下赋役、赈济与赏赐的记录。\footnote{核心来源：《三国志·蜀书·后主传》。材料限度：本纪可证实所载命令、任免或结果；不据此推定各郡执行、全部参与者或私人动机。；本纪选择朝廷可见事项；不得由一条记录推定各郡社会状况或具体群体动机。；本纪可核对行动的基本顺序；兵额、路线、战损和具体指挥过程仍须另证。}
\begin{itemize}[leftmargin=1.8em,itemsep=2pt,topsep=3pt]
  \item \eventanchor{TK-SRC-0232-03}{十二年春二月，亮由斜谷出，始以流馬運}。
  \item \eventanchor{TK-SRC-0232-04}{十年，亮休士勸農於黃沙，作流馬木牛畢，教兵講武}。
  \item \eventanchor{TK-SRC-0232-05}{十四年夏四月，後主至湔，登觀坂，看汶水之流，旬日還成都}。
  \item \eventanchor{TK-SRC-0232-06}{十五年夏六月，皇后張氏薨}。
  \item \eventanchor{TK-SRC-0232-07}{徙武都氐王苻健及氐民四百餘戶於成都}。
  \item \eventanchor{TK-SRC-0232-08}{夏四月，進蔣琬位為大將軍}。
  \item \eventanchor{TK-SRC-0232-09}{以丞相留府長史蔣琬為尚書令，總統國事}。
  \item \eventanchor{TK-SRC-0232-10}{征西大將軍魏延與丞相長史楊儀爭權不和，舉兵相攻，延敗走}。
\end{itemize}

\paragraph{238年} 《三国志·蜀书·后主传》在这一年并列留下朝廷可见的处置与记述的记录。\footnote{核心来源：《三国志·蜀书·后主传》。材料限度：本纪选择朝廷可见事项；不得由一条记录推定各郡社会状况或具体群体动机。}
\begin{itemize}[leftmargin=1.8em,itemsep=2pt,topsep=3pt]
  \item \eventanchor{TK-SRC-0238-02}{立子睿為太子，子瑤為安定王}。
  \item \eventanchor{TK-SRC-0238-03}{延熙元年春正月，立皇后張氏}。
\end{itemize}

\paragraph{239年} 《三国志·蜀书·后主传》在这一年并列留下朝廷可见的处置与记述的记录。\footnote{核心来源：《三国志·蜀书·后主传》。材料限度：本纪可证实所载命令、任免或结果；不据此推定各郡执行、全部参与者或私人动机。}
\begin{itemize}[leftmargin=1.8em,itemsep=2pt,topsep=3pt]
  \item \eventanchor{TK-SRC-0239-01}{二年春三月，進蔣琬位為大司馬}。
\end{itemize}

\paragraph{240年} 《三国志·蜀书·后主传》在这一年并列留下官署、任命与诏令的记录。\footnote{核心来源：《三国志·蜀书·后主传》。材料限度：本纪可核对行动的基本顺序；兵额、路线、战损和具体指挥过程仍须另证。}
\begin{itemize}[leftmargin=1.8em,itemsep=2pt,topsep=3pt]
  \item \eventanchor{TK-SRC-0240-01}{三年春，使越巂太守張嶷平定越巂郡}。
\end{itemize}

\paragraph{247年} 《三国志·蜀书·后主传》在这一年并列留下官署、任命与诏令的记录。\footnote{核心来源：《三国志·蜀书·后主传》。材料限度：本纪选择朝廷可见事项；不得由一条记录推定各郡社会状况或具体群体动机。；本纪可核对行动的基本顺序；兵额、路线、战损和具体指挥过程仍须另证。；本纪可证实所载命令、任免或结果；不据此推定各郡执行、全部参与者或私人动机。}
\begin{itemize}[leftmargin=1.8em,itemsep=2pt,topsep=3pt]
  \item \eventanchor{TK-SRC-0247-01}{冬，拔狄道、河間[河關]、臨洮三縣民，居於綿竹、繁縣}。
  \item \eventanchor{TK-SRC-0247-02}{經退保狄道城，維卻住鍾題}。
  \item \eventanchor{TK-SRC-0247-03}{秋，涪陵屬國民夷反，車騎將軍鄧芝往討，皆破平之}。
  \item \eventanchor{TK-SRC-0247-04}{秋，衛將軍姜維出攻雍州，不克而還}。
  \item \eventanchor{TK-SRC-0247-05}{秋八月，維為魏大將軍鄧艾所破於上邽}。
  \item \eventanchor{TK-SRC-0247-06}{十二年春正月，魏誅大將軍曹爽等，右將軍夏侯霸來降}。
  \item \eventanchor{TK-SRC-0247-07}{十六年春正月，大將軍費禕為魏降人郭循所殺於漢壽}。
  \item \eventanchor{TK-SRC-0247-08}{十七年春正月，姜維還成都}。
  \item \eventanchor{TK-SRC-0247-09}{十三年，姜維復出西平，不克而還}。
  \item \eventanchor{TK-SRC-0247-10}{十四年夏，大將軍費禕還成都}。
\end{itemize}

\paragraph{258年} 《三国志·蜀书·后主传》在这一年并列留下官署、任命与诏令的记录。\footnote{核心来源：《三国志·蜀书·后主传》。材料限度：本纪选择朝廷可见事项；不得由一条记录推定各郡社会状况或具体群体动机。}
\begin{itemize}[leftmargin=1.8em,itemsep=2pt,topsep=3pt]
  \item \eventanchor{TK-SRC-0258-03}{史官言景星見，於是大赦，改年}。
\end{itemize}

\paragraph{259年} 《三国志·蜀书·后主传》在这一年并列留下朝廷可见的处置与记述的记录。\footnote{核心来源：《三国志·蜀书·后主传》。材料限度：本纪选择朝廷可见事项；不得由一条记录推定各郡社会状况或具体群体动机。}
\begin{itemize}[leftmargin=1.8em,itemsep=2pt,topsep=3pt]
  \item \eventanchor{TK-SRC-0259-01}{二年夏六月，立子諶為北地王，恂為新興王，虔為上黨王}。
\end{itemize}

\paragraph{260年} 《三国志·蜀书·后主传》在这一年并列留下官署、任命与诏令的记录。\footnote{核心来源：《三国志·蜀书·后主传》。材料限度：本纪可核对行动的基本顺序；兵额、路线、战损和具体指挥过程仍须另证。}
\begin{itemize}[leftmargin=1.8em,itemsep=2pt,topsep=3pt]
  \item \eventanchor{TK-SRC-0260-02}{三年秋九月，追諡故將軍關羽、張飛、馬超、龐統、黃忠}。
\end{itemize}

\paragraph{261年} 《三国志·蜀书·后主传》在这一年并列留下官署、任命与诏令的记录。\footnote{核心来源：《三国志·蜀书·后主传》。材料限度：本纪可核对行动的基本顺序；兵额、路线、战损和具体指挥过程仍须另证。}
\begin{itemize}[leftmargin=1.8em,itemsep=2pt,topsep=3pt]
  \item \eventanchor{TK-SRC-0261-01}{四年春三月，追諡故將軍趙雲}。
\end{itemize}

\paragraph{263年} 《三国志·蜀书·后主传》在这一年并列留下朝廷可见的处置与记述的记录。\footnote{核心来源：《三国志·蜀书·后主传》。材料限度：本纪可证实所载命令、任免或结果；不据此推定各郡执行、全部参与者或私人动机。}
\begin{itemize}[leftmargin=1.8em,itemsep=2pt,topsep=3pt]
  \item \eventanchor{TK-SRC-0263-01}{艾得書，大喜，即報書，遣紹良先還}。
  \item \eventanchor{TK-SRC-0263-02}{艾使後主止其故宮，身往造焉}。
\end{itemize}

\paragraph{264年} 《三国志·蜀书·后主传》在这一年并列留下赋役、赈济与赏赐的记录。\footnote{核心来源：《三国志·蜀书·后主传》。材料限度：本纪所载多为朝廷命令或结果，不足以证明地方执行、总量和受影响户数。}
\begin{itemize}[leftmargin=1.8em,itemsep=2pt,topsep=3pt]
  \item \eventanchor{TK-SRC-0264-01}{食邑萬戶，賜絹萬匹，奴婢百人，他物稱是}。
\end{itemize}
